\section{A caixa}

\textbf{(1.a)} Mostre que, na fase de contato à direita, a equação de movimento é
\begin{center}
    $m \ddot{\vec{x}} = -k A_R \vec{x} + k_{muro} d \hat{x_N}, \quad  A_R = A_0 + \frac{k_{muro}}{k} \hat{x_N} \hat{x_N}^T$,
\end{center}
onde $x_N {x_N}^T$ é a matriz cujo único elemento não nulo é o $(N, N)$, igual a 1. O termo constante
não existia no Projeto 1, porque lá a parede estava na posição de equilíbrio da massa $N$. Mostre
que a mudança de variáveis $\vec{y} = \vec{x} - \vec{u}$, com $A_r \vec{u} = \frac{k_{muro}}{k} d \hat{x_N}$, elimina esse termo e devolve
a forma do Projeto 1, $m \ddot{\vec{y}} = -k A_R \vec{y}$. Encontre $\vec{u}$ explicitamente e interprete-o fisicamente. Faça
o mesmo para o contato à esquerda.

\vspace{0.3cm}
\textbf{Solução:} Assim como  projeto 1, na configuração livre, a equação de movimento é dado por:
\begin{center}
    $m \ddot{\vec{x}} = -k A_0 \vec{x}$ 
\end{center}


Quando a mola está em contato com a parede direita, significa que a massa $N$ sofre uma força a mais da parede. Essa força é dado por:

\begin{center}
    $F_N = -k_{muro} (x_N - d)$ 
\end{center}

Lembre da força que calculamos no projeto 1 para a massa $N$:
\begin{center}
    $m \ddot{x}_N = -k (x_N - x_{N-1}) $
\end{center}

Com a força adicional da parede, temos que a equação de movimento da massa $N$ é:
\begin{center}
    $m \ddot{x}_N = -k (x_N - x_{N-1}) - k_{muro} (x_N - d)$
\end{center}

Que pode ser escrito como:

\begin{center}
    $m \ddot{x}_N = -k [-x_{N-1}+(1+\frac{k_{muro}}{k}) x_N] + k_{muro} d$ 
\end{center}

Perceba que no contato à direita, somente a massa $N$ sofre uma alteração na força. As outras massas de $1$ até $N-1$ não sofrem alteração.
Podemos escrever isso na forma matricial da seguinte forma:
\begin{center}
    $m \ddot{\vec{x}} = -k A_0 \vec{x} -k_{muro} \hat{x}_N \hat{x}_N^T \vec{x}+k_{muro} d \hat{x_N}$
\end{center}
onde $\hat{x}_N = (0,0, \dots, 1)$.
Os termos $-k_{muro} \hat{x}_N {\hat{x}_N}^T \vec{x}$ e $k_{muro} d \hat{x_N}$ só alteram $m \ddot{x}_N$ na matriz, são as forças adicionais que a parede faz na massa $N$.

Vamos ficar com

\begin{center}
    $m \ddot{\vec{x}} = -k (A_0 \vec{x}+\frac{k_{muro}}{k} \hat{x_N} \hat{x}_N ^ T) \vec{x} + k_{muro} d \hat{x}_N$
\end{center}
Se definimos $A_R = A_0 \vec{x}+\frac{k_{muro}}{k} \hat{x_N} \hat{x}_N ^ T$, teremos justamente o que queremos mostrar.
\begin{center}
    $m \ddot{\vec{x}} = -k A_R \vec{x} + k_{muro} d \hat{x}_N$
\end{center}

Colocando $\vec{y} = \vec{x} - \vec{u}$ com $A_R \vec{u} = (k_{muro}/k) d \hat{x}_N$. O vetor $\vec{u}$ é constante. 
Por causa disso, vamos ter que $\ddot{\vec{y}} = \ddot{\vec{x}}$.
E então, isolando $\vec{x}$, $\vec{x} = \vec{y} + \vec{u}$ 

\begin{align*}
    m \ddot{\vec{y}} = m \ddot{\vec{x}} &= -k A_R (\vec{y}+\vec{u}) +k_{muro} d \hat{x}_N \\
    &= -k A_R \vec{y} -k A_R \vec{u} +k_{muro} d \hat{x}_N \\
    &=-k A_R \vec{y} - k_{muro} d \hat{x}_N + k_{muro} d \hat{x}_N = -k A_R \vec{y}
\end{align*}

Queremos encontrar $\vec{u}$ tal que
\begin{center}
    $A_R \vec{u} = (k_{muro}/k) d \hat{x}_N$
\end{center}

\[
    A_R \vec{u} =
    \begin{pmatrix}
        1 & -1 & 0 & \cdots & 0 \\
        -1 & 2 & -1 & \cdots & 0 \\
        0 & -1 & \ddots & \ddots & \vdots \\
        \vdots & \vdots & \ddots & 2 & -1 \\
        0 & 0 & \cdots & -1 & 1 + \frac{k_{\text{muro}}}{k}
    \end{pmatrix}
    \begin{pmatrix}
        u_1 \\
        u_2 \\
        \vdots \\
        u_{N-1} \\
        u_{N}
    \end{pmatrix}
    =
    \begin{pmatrix}
        0 \\
        0 \\
        \vdots \\
        0 \\
        \frac{k_{\text{muro}}}{k} d
    \end{pmatrix}
\]

Isso é um sistema de equações que vai nos dar que:

\begin{align*}
    u_1 - u_2 &= 0 \\
    -u_1 + 2u_2 - u_3 &= 0 \\
    &\vdots \\
    -u_{N-3} + 2u_{N-2} - u_{N-1} &= 0 \\
    -u_{N-1} + \left(1 + \frac{k_{muro}}{k}\right)u_N &= \frac{k_{muro}}{k} d
\end{align*}

É fácil ver que $u_1 = u_2$, usando isso na segunda equação nos dá que $u_1-u_3=0$. E então $u_1=u_2=u_3$.
Com as primeiras $N-1$ equações, podemos perceber que $u_1 = u_2 = \dots = u_{N-1} = u_N$. Usando isso na última equação, vamos ter que
\begin{center}
    $-u_{N}+(1+\frac{k_{muro}}{k}) u_N = \frac{k_{muro}}{k} d$
\end{center}


Isso vai nos dá que $u_{N} = d$. Portanto $\vec{u} = (d, d, \dots, d)$
O vetor $\vec{u}$ pode ser interpretado como um vetor constante que translada todas as massas.
Para o contato com a parede esquerda. Teremos uma força adicional na massa $1$ e montamos a matriz de forma análoga.

\vspace{0.3cm}
\textbf{(1.b)} Diga por que $A_R$ não tem autovalor nulo, ao contrário de $A_0$. Adapte a rotina do item
\textbf{(5.a)} do Projeto 1 para que ela receba, além da matriz e das condições iniciais, o deslocamento $\vec{u}$
do ponto de equilíbrio, e devolva $\vec{x}(t)$ e $\vec{v}(t)$ em qualquer fase.

\vspace{0.3cm}
\textbf{Solução:}
O determinante da matriz $A_R$ é não nulo. Se uma matriz tem determinante não nulo, ela não pode ter autovalor nulo.
$A_0$ tem autovalor nulo, como vimos no projeto 1. Isso implica que o seu determinante é nulo.
Podemos ver que o determinante não é nulo, calculando ele por meio do método de chió. Reduzindo o tamanho da matriz e obtendo uma recorrencia para os determinantes.
O código adaptado está no main.ipynb.


\vspace{0.3cm}

\textbf{(1.c)} A simulação é uma sequência de fases. Em cada uma, a condição inicial é o estado $(\vec{x},\vec{v})$
no fim da anterior, e o fim da fase é o primeiro instante em que uma das condições de contato
muda: $x_N = d$ ou $x_1 = -d$ na fase livre, e o descolamento na fase de contato. Como $\vec{x}(t)$ é
conhecido em forma fechada, achar esse instante é achar a primeira raiz de uma função conhecida.
Proponha e implemente um método para isso. Uma possibilidade é amostrar a função numa
grade de passo $\Delta t$, achar a primeira troca de sinal e refinar a raiz por bisseção. Nesse caso,
explique como escolher $\Delta t$, e o que pode dar errado se ele for grande demais. Estime a duração
típica de um contato.

\vspace{0.3cm}
\textbf{Solução:}
Queremos determinar a primeira raiz para a função $f(t) = x_N(t)-d = 0$(ou $f(t) = x_1(t)+d = 0$, depende da fase). Para isso podemos discretizar o tempo $\delta t$ e a cada passo tentar avaliar se houve troca de sinal na função.
Temos uma fórmula que estima o tempo de contato:
\begin{center}
    $t_{contato} \approx \pi \sqrt{\frac{m}{k_{muro}}}$
\end{center}

Para ver essa fórmula, podemos pensar que o período completo de um ciclo é $T = 2 \pi \sqrt{\frac{m}{k_{muro}}}$
e ela gasta meio período $(T/2)$ em contato com a parede

\vspace{2.46cm}

\textbf{(1.d)} Escreva a energia mecânica total, incluindo as molas das paredes quando comprimidas, e
verifique que ela se conserva ao longo de milhares de colisões. Explique como o erro na energia
depende da precisão com que os instantes de troca de fase são calculados.

\vspace{0.3cm}
\textbf{Solução:}
A energia mecânica é a soma da energia cinética
$K = \frac{1}{2} m \sum_{i=1}^{N} \ddot{x}_i$ da energia potencial interna $U_{interna} = \frac{1}{2} k \sum_{i=1}^{N-1} (x_{i+1}-x_{i})^2$ e da energia potencial das paredes.
Na fase livre, essa energia das paredes será zero. Em contato com a parede direita, a energia potencial dessa parede será $U_{muro} = \frac{1}{2} k_{muro} (x_N-d)^2$ (somente a mola da parede direita deformou, ela vai ter energia potencial para voltar a sua posição relaxada). 
Para o contato com a parede esquerda, a energia potencial é $U_{muro} = \frac{1}{2} k_{muro} (x_1+d)^2$.
Essa energia vai se conservar mesmo após milhares de colisões. Nosso sistema não tem atrito (nem outras forças dissipativas). Tanto a energia potencial das molas quanto das paredes se derivam de potenciais elásticos.

\vspace{0.3cm}
\textbf{(1.e)} Como no Projeto 1, é obrigatória uma animação das massas em movimento, agora com
as duas paredes. Além dela, produzam as visualizações que acharem que revelam melhor o
que está acontecendo. As sugestões do item \textbf{(5.d)} do Projeto 1 continuam valendo; o diagrama
espaço-tempo, com as paredes como duas retas verticais, é particularmente instrutivo. O critério
é o mesmo: a figura deve tornar evidente alguma coisa que não era evidente antes.
Como no Projeto 1, o padrão é $k_{muro} = k$. Como ponto de partida, tomem $N = 10$ e $d = 6$, com
a condição inicial do Projeto 1,
\begin{center}
    $\vec{x}^0 = (0, 0, \dots , 0)$,  $\vec{v}^0 = (v, v, \dots , v)$, $v = 1$.
\end{center}

Para $k_{muro} \gg k$, a massa da ponta pode tocar a parede várias vezes seguidas, como vocês viram
no Projeto 1. Chamaremos de colisão o conjunto de todos os contatos com uma mesma parede
entre duas visitas à parede oposta.

\vspace{0.3cm}
\textbf{Solução:}